\documentclass[UTF8,a4paper,10pt]{ctexart}
\usepackage{ipara-handbook}
\usepackage{booktabs,tabularx,tikz}
\usetikzlibrary{arrows.meta,positioning}
\begin{document}
\section*{H-I01\quad 微积分建模：从局部微元到整体量}
\textbf{定位：}本 Integration 只验收跨 Topic 组合，不重新教授单个微元或积分机制。
\begin{tcolorbox}[title=母流程]
目标量 $\longrightarrow$ 选对象与坐标 $\longrightarrow$ 切出局部元
$\longrightarrow$ 写 $dQ$ 与积分域 $\longrightarrow$ 累积
$\longrightarrow$ 单位/边界/方向复核。
\end{tcolorbox}
\begin{center}
\begin{tikzpicture}[node distance=7mm, every node/.style={draw,rounded corners,align=center,minimum width=2.8cm,minimum height=8mm,font=\small},>=Latex]
\node (a) {目标量\\面积/质量/功};
\node (b) [right=of a] {局部元\\$dQ$};
\node (c) [right=of b] {区域或路径\\$\Omega,\ C,\ S$};
\node (d) [right=of c] {积分与复核};
\draw[->,thick] (a)--(b); \draw[->,thick] (b)--(c); \draw[->,thick] (c)--(d);
\end{tikzpicture}
\end{center}
\subsection*{1.\ 平面薄片质量：密度先进入微元}
对区域 $\Omega$ 上的面密度 $\rho(x,y)$，总质量先写
\[
dM=\rho(x,y)\,dA,\qquad M=\iint_{\Omega}\rho(x,y)\,dA.
\]
若 $\Omega$ 关于 $y$ 轴对称且 $\rho$ 为奇函数 $x$，则带符号积分发生抵消；物理质量仍要求 $\rho\ge0$，必须先检查题设是否允许该密度。
\subsection*{2.\ 体积、弧长与功的切换}
\begin{tabularx}{\linewidth}{@{}lX@{}}
\toprule
目标 & 典型微元与资格\\ \midrule
旋转体体积 & $dV=\pi r^2\,dx$ 或壳元 $dV=2\pi r h\,dx$；先确认切片方向\\
曲线弧长 & $ds=\sqrt{(dx)^2+(dy)^2}$；参数式用 $\sqrt{(x')^2+(y')^2}\,dt$\\
力做功 & $dW=\mathbf F\cdot d\mathbf r$；方向和单位由位移元决定\\
曲面通量 & $d\Phi=\mathbf F\cdot\mathbf n\,dS$；必须先定向\\
\bottomrule
\end{tabularx}
\subsection*{3.\ 旋转体：选择的是切片方向，不是背公式}
同一平面区域绕轴旋转时有两种局部化：垂直于轴切片得到圆盘/垫圈，平行于轴切片得到圆柱壳。
若已知 $y=f(x)$ 且绕 $y$ 轴，比较两条路线：反解 $x=g(y)$ 后用圆盘，或保留 $x$ 用壳元
\[
dV=2\pi\,\underbrace{|x-x_{\rm axis}|}_{\text{半径}}
\underbrace{h(x)}_{\text{壳高}}\,dx.
\]
选择标准是区域描述、反解成本与\textbf{旋转后的覆盖次数}；无论哪条路线，都必须保证切片不重不漏。

\begin{warnbox}{壳层真正按“半径”编号，而不是按原坐标编号}
若区域跨过旋转轴，轴两侧的 $x=x_{\rm axis}\pm r$ 会旋转成\textbf{同一个半径 $r$ 的圆柱壳}。因此不能机械把
\[
2\pi\int |x-x_{\rm axis}|h(x)\,dx
\]
在轴两侧全部相加，否则可能重复覆盖。应改用半径 $r\ge0$：先把同一 $r$ 在两侧产生的高度区间取并集，记其总长度为 $H(r)$，再积
\[
\boxed{V=2\pi\int rH(r)\,dr.}
\]
最小反例是矩形 $[-1,1]\times[0,1]$ 绕 $y$ 轴：实体只是体积 $\pi$ 的单位圆柱；把左右壳直接相加会错误得到 $2\pi$。所以壳层法的第一问不是“半径是多少”，而是“\textbf{这个半径是否已经被另一侧生成过？}”
\end{warnbox}

圆盘/垫圈也应按“半径集合”理解：一个垂直于轴的线段若跨过旋转轴，旋转后内半径是 $0$，形成圆盘；不能把上下边界到轴的距离机械做平方差。

若典型连通平面区域 $D$ 绕同一平面内、且不穿过 $D$（因而 $D$ 位于轴的一侧）的直线
\[
L:ax+by+c=0
\]
旋转一周，并且面积元扫出的体积按一次覆盖计数，可以把每个面积元 $dA$ 看成绕轴扫出的薄环体：
\[
dV=2\pi r(x,y)\,dA,
\qquad
r(x,y)=\frac{|ax+by+c|}{\sqrt{a^2+b^2}}.
\]
于是
\[
\boxed{V=\iint_D2\pi r(x,y)\,dA.}
\]
这不是新的孤立公式，而是壳层思想从“一条竖条”推广到“任意面积元”。“轴不穿过区域 + 一次覆盖”是不可省略的资格；若区域分离在轴两侧、存在自重叠或旋转后不同面积元生成同一体积层，仍必须回到覆盖审计，不能机械累加。若满足上述资格，再利用质心定义可压成 $V=2\pi A\,d(\bar P,L)$。

旋转曲面面积使用另一种微元：
\[
dS=2\pi\,\underbrace{r}_{\text{到旋转轴距离}}\,ds.
\]
它局部累加“周长 $\times$ 弧长”，不能把 $ds$ 静默替换为 $dx$。若生成曲线存在重复参数化、自交或跨轴回折，还要检查不同微段旋转后是否扫到同一块曲面；否则积分会按生成重数计面积。体积有限也不推出表面积有限，两者尾部模型不同。

\subsection*{4.\ 相关变化率：保持动态对象直到求导完成}
把几何关系写成时间的恒等式 $G(x(t),y(t),\ldots)=0$，先对 $t$ 求导：
\[
G_x\frac{dx}{dt}+G_y\frac{dy}{dt}+\cdots=0,
\]
再代入题目给定时刻的长度、半径或体积。若求导前就把瞬时值代成常数，对应变化率会被错误地消掉。
单位是这里的接口语言：$\mathrm m/\mathrm s$、$\mathrm m^2/\mathrm s$、$\mathrm m^3/\mathrm s$ 分别提示长度、面积、体积的变化率对象。

\subsection*{5.\ 压力与做功：强度必须乘对几何载体}
变力做功沿位移累积 $dW=F(x)\,dx$；抽水时局部贡献为
\[
dW=\underbrace{\rho g A(y)\,dy}_{\text{薄层重量}}
\underbrace{D(y)}_{\text{提升距离}}.
\]

\begin{tcolorbox}[title=质心法为什么能压缩整段抽水积分]
若所有质量元最终都被提升到同一个出口高度 $H$，且只计算克服重力所做的功，则
\[
W=\int g(H-y)\,dm
=Mg\left(H-\frac1M\int y\,dm\right)
=Mg(H-\bar y).
\]
所以“把全部质量集中在质心”不是新的经验公式，而是\textbf{积分线性性把所有高度差的质量加权平均压缩成质心高度}。直觉上，每一小块都要爬不同距离；质心记录的正是这些距离的质量加权平均。

这条压缩只有在最终高度统一、目标量只有重力势能变化时才成立。若不同质量元终点不同，或还要计入压强损失、其他能量项，就必须回到逐层微元，不能套质心式。
\end{tcolorbox}

静水压力则在面积上累积，取从液面量起的竖直深度 $h$：
\[
dF=\rho gh\,dA=\rho gh\,L(h)\,dh.
\]
“深度”不是任意坐标值，“移动距离”也不是切片位置；二者都要从物理对象重新定义。

\subsection*{6.\ 最小组合例：变密度细杆}
细杆位于 $x\in[a,b]$，线密度 $\lambda(x)$，则
\[
dM=\lambda(x)\,dx,\quad M=\int_a^b\lambda(x)\,dx,\quad
\bar x=\frac{1}{M}\int_a^b x\lambda(x)\,dx.
\]
先求总量 $M$ 才能定义质心；若 $\lambda$ 在区间变号，需先判断其物理意义。
\begin{tcolorbox}[title=验收检查]
微元单位是否与目标量一致？区域/路径/曲面边界是否完整？
对称性作用于完整 integrand 还是只作用于几何域？有向积分是否明确方向？
\end{tcolorbox}
\subsection*{7.\ Anti-Bridge}
单位检查只能发现量纲冲突，不能证明切片、区域或方向正确；“写出了积分”也不等于完成建模。若局部元无法从对象和坐标唯一解释，应回退到 Topic01/05/08/09 Owner。
\end{document}
